% GENERATED FILE. Edit canonical lessons, then run npm run generate:books.
% canonical-source-hash: fnv1a64:9ecbb0211dc1821d
% canonical-lessons: MR-C296-dhadakne, MR-C296-adalne, MR-C296-besuddha, MR-C296-savarne, MR-C296-hukne

\chapter{A Crash on the Road}
\label{ch:mr-a-crash-on-the-road}
\hlchaptermodality{296}

\begin{chapteropening}
\textbf{By the end of this chapter:} \emph{I can describe a road accident and how someone came through it.}

Complete the last lesson of chapter 296: I can describe a road accident and how someone came through it.
\end{chapteropening}

\section[\texorpdfstring{\mr{धडकणे} (dhaḍakṇe)}{धडकणे (dhaḍakṇe)}]{\mr{धडकणे} (dhaḍakṇe) — to collide; to crash}

\label{lesson:MR-C296-dhadakne}



\begin{quote}
Before the new one: say the Marathi for a crisis, then the Marathi for dreadful.
\end{quote}

\subsection*{You'll want to know: \mr{धडकणे}}
\textbf{\mr{धडकणे}} — \emph{dhaḍakṇe} — \textquotedblleft{}to collide; to crash\textquotedblright{}.

\begin{quote}
\mr{भरधाव} \mr{गाडी} \mr{झाडावर} \mr{धडकली}.
\end{quote}

\begin{quote}
\emph{bhardhāv gāḍī jhāḍāvar dhaḍaklī.}
\end{quote}

\begin{quote}
The speeding car crashed into a tree.
\end{quote}

Past (he …): \mr{तो} \mr{धडकला}, \emph{to dhaḍaklā}.

What is hit takes -\mr{वर} or -\mr{ला}: \mr{भिंतीवर} \mr{धडकणे}, to crash into a wall.

\begin{cousinweb}
\textbf{\mr{धडकणे}} (\emph{dhaḍakṇe}) is inherited: CDIAL lists Marathi \emph{dhaḍakṇe}, 'to throb', under \emph{*dhaḍ}, 'throb, beat' (CDIAL 6711); the sense of a crash grows from the thud.
\end{cousinweb}

\begin{grammarlens}[title={What you've built}]
One more word for telling what happened.
\end{grammarlens}

\subsection*{Your turn}
\begin{itemize}
\raggedright
  \item \emph{Say it:} \emph{dhaḍakṇe}
  \item \emph{Say it:} \emph{dhaḍakṇe}, once more
  \item \emph{Say it:} \emph{bhardhāv gāḍī jhāḍāvar dhaḍaklī.} — The speeding car crashed into a tree.
  \item \emph{Recall:} say the Marathi for a crisis, then the Marathi for dreadful, then say \emph{dhaḍakṇe} again
\end{itemize}

\begin{tcolorbox}[breakable,colback=teal!4,colframe=teal!35!black,title={Before you move on}]
What does \mr{धडकणे} mean? (\textquotedblleft{}To collide; to crash\textquotedblright{}.) Say it once more.
\end{tcolorbox}
\section[\texorpdfstring{\mr{आदळणे} (ādaḷṇe)}{आदळणे (ādaḷṇe)}]{\mr{आदळणे} (ādaḷṇe) — to slam into; to dash against}

\label{lesson:MR-C296-adalne}



\begin{quote}
Before the new one: say the Marathi for dreadful, then the Marathi for to collide.
\end{quote}

\subsection*{You'll want to know: \mr{आदळणे}}
\textbf{\mr{आदळणे}} — \emph{ādaḷṇe} — \textquotedblleft{}to slam into; to dash against\textquotedblright{}.

\begin{quote}
\mr{सायकल} \mr{घसरली} \mr{आणि} \mr{तो} \mr{खांबावर} \mr{आदळला}.
\end{quote}

\begin{quote}
\emph{sāykal ghasarlī āṇi to khāmbāvar ādaḷlā.}
\end{quote}

\begin{quote}
The bicycle skidded and he slammed into a pole.
\end{quote}

Past (he …): \mr{तो} \mr{आदळला}, \emph{to ādaḷlā}.

The thing struck takes -\mr{वर}: \mr{जमिनीवर} \mr{आदळणे}, to hit the ground hard.

\begin{cousinweb}
\textbf{\mr{आदळणे}} (\emph{ādaḷṇe}) is inherited: CDIAL lists Marathi \emph{adaḷṇẽ}, 'to dash down or against; to rush against', beside \emph{ādaḷṇẽ}, under \emph{*ādalati}, 'crushes' (CDIAL 1144).
\end{cousinweb}

\begin{grammarlens}[title={What you've built}]
One more word for telling what happened.
\end{grammarlens}

\subsection*{Your turn}
\begin{itemize}
\raggedright
  \item \emph{Say it:} \emph{ādaḷṇe}
  \item \emph{Say it:} \emph{ādaḷṇe}, once more
  \item \emph{Say it:} \emph{sāykal ghasarlī āṇi to khāmbāvar ādaḷlā.} — The bicycle skidded and he slammed into a pole.
  \item \emph{Recall:} say the Marathi for dreadful, then the Marathi for to collide, then say \emph{ādaḷṇe} again
\end{itemize}

\begin{tcolorbox}[breakable,colback=teal!4,colframe=teal!35!black,title={Before you move on}]
What does \mr{आदळणे} mean? (\textquotedblleft{}To slam into; to dash against\textquotedblright{}.) Say it once more.
\end{tcolorbox}
\section[\texorpdfstring{\mr{बेशुद्ध} (beśuddha)}{बेशुद्ध (beśuddha)}]{\mr{बेशुद्ध} (beśuddha) — unconscious}

\label{lesson:MR-C296-besuddha}



\begin{quote}
Before the new one: say the Marathi for to collide, then the Marathi for to slam into.
\end{quote}

\subsection*{You'll want to know: \mr{बेशुद्ध}}
\textbf{\mr{बेशुद्ध}} — \emph{beśuddha} — \textquotedblleft{}unconscious\textquotedblright{}.

\begin{quote}
\mr{अपघातानंतर} \mr{तो} \mr{काही} \mr{मिनिटं} \mr{बेशुद्ध} \mr{होता}.
\end{quote}

\begin{quote}
\emph{apghātānantar to kāhī miniṭaṁ beśuddha hotā.}
\end{quote}

\begin{quote}
After the accident he was unconscious for a few minutes.
\end{quote}

\mr{बेशुद्ध} \mr{होणे}, to faint; \mr{शुद्धीवर} \mr{येणे}, to come round.

\begin{cousinweb}
\textbf{\mr{बेशुद्ध}} (\emph{beśuddha}) is built from the Persian prefix \emph{be-}, 'without', and Sanskrit \emph{śuddha}, 'clear, pure' (Monier-Williams), here in the sense of a clear mind.
\end{cousinweb}

\begin{grammarlens}[title={What you've built}]
One more word for telling what happened.
\end{grammarlens}

\subsection*{Your turn}
\begin{itemize}
\raggedright
  \item \emph{Say it:} \emph{beśuddha}
  \item \emph{Say it:} \emph{beśuddha}, once more
  \item \emph{Say it:} \emph{apghātānantar to kāhī miniṭaṁ beśuddha hotā.} — After the accident he was unconscious for a few minutes.
  \item \emph{Recall:} say the Marathi for to collide, then the Marathi for to slam into, then say \emph{beśuddha} again
\end{itemize}

\begin{tcolorbox}[breakable,colback=teal!4,colframe=teal!35!black,title={Before you move on}]
What does \mr{बेशुद्ध} mean? (\textquotedblleft{}Unconscious\textquotedblright{}.) Say it once more.
\end{tcolorbox}
\section[\texorpdfstring{\mr{सावरणे} (sāvarṇe)}{सावरणे (sāvarṇe)}]{\mr{सावरणे} (sāvarṇe) — to recover; to steady oneself}

\label{lesson:MR-C296-savarne}



\begin{quote}
Before the new one: say the Marathi for to slam into, then the Marathi for unconscious.
\end{quote}

\subsection*{You'll want to know: \mr{सावरणे}}
\textbf{\mr{सावरणे}} — \emph{sāvarṇe} — \textquotedblleft{}to recover; to steady oneself\textquotedblright{}.

\begin{quote}
\mr{थोड्या} \mr{वेळाने} \mr{तो} \mr{त्या} \mr{भीतीतून} \mr{सावरला}.
\end{quote}

\begin{quote}
\emph{thoḍyā veḷāne to tyā bhītītūn sāvarlā.}
\end{quote}

\begin{quote}
After a while he recovered from that fright.
\end{quote}

Past (he …): \mr{तो} \mr{सावरला}, \emph{to sāvarlā}.

The trouble takes -\mr{तून}: \mr{भीतीतून} \mr{सावरणे}. Transitive too: \mr{स्वतःला} \mr{सावरणे}, to steady oneself.

\begin{cousinweb}
\textbf{\mr{सावरणे}} (\emph{sāvarṇe}) is inherited: CDIAL lists Marathi \emph{sāvarṇe}, 'to gather up, recover, rectify', under \emph{*saṁvarati}, 'restrains' (CDIAL 13015).
\end{cousinweb}

\begin{grammarlens}[title={What you've built}]
One more word for telling what happened.
\end{grammarlens}

\subsection*{Your turn}
\begin{itemize}
\raggedright
  \item \emph{Say it:} \emph{sāvarṇe}
  \item \emph{Say it:} \emph{sāvarṇe}, once more
  \item \emph{Say it:} \emph{thoḍyā veḷāne to tyā bhītītūn sāvarlā.} — After a while he recovered from that fright.
  \item \emph{Recall:} say the Marathi for to slam into, then the Marathi for unconscious, then say \emph{sāvarṇe} again
\end{itemize}

\begin{tcolorbox}[breakable,colback=teal!4,colframe=teal!35!black,title={Before you move on}]
What does \mr{सावरणे} mean? (\textquotedblleft{}To recover; to steady oneself\textquotedblright{}.) Say it once more.
\end{tcolorbox}
\section[\texorpdfstring{\mr{हुकणे} (hukṇe)}{हुकणे (hukṇe)}]{\mr{हुकणे} (hukṇe) — to be missed; to miss a chance}

\label{lesson:MR-C296-hukne}



\begin{quote}
Before the new one: say the Marathi for unconscious, then the Marathi for to recover.
\end{quote}

\subsection*{You'll want to know: \mr{हुकणे}}
\textbf{\mr{हुकणे}} — \emph{hukṇe} — \textquotedblleft{}to be missed; to miss a chance\textquotedblright{}.

\begin{quote}
\mr{मला} \mr{उशीर} \mr{झाला} \mr{आणि} \mr{माझी} \mr{बस} \mr{हुकली}.
\end{quote}

\begin{quote}
\emph{malā uśīr jhālā āṇi mājhī bas huklī.}
\end{quote}

\begin{quote}
I was late and missed my bus.
\end{quote}

Past (he …): \mr{तो} \mr{हुकला}, \emph{to huklā}.

As with \mr{चुकणे}, the thing missed is the subject and agrees: \mr{बस} \mr{हुकली}, \mr{संधी} \mr{हुकली}, the chance was missed.

\begin{cousinweb}
\textbf{\mr{हुकणे}} (\emph{hukṇe}) is a Marathi verb for a chance, a bus or an aim that slips by; its origin is not settled.
\end{cousinweb}

\begin{grammarlens}[title={What you've built}]
One more word for telling what happened.
\end{grammarlens}

\subsection*{Your turn}
\begin{itemize}
\raggedright
  \item \emph{Say it:} \emph{hukṇe}
  \item \emph{Say it:} \emph{hukṇe}, once more
  \item \emph{Say it:} \emph{malā uśīr jhālā āṇi mājhī bas huklī.} — I was late and missed my bus.
  \item \emph{Recall:} say the Marathi for unconscious, then the Marathi for to recover, then say \emph{hukṇe} again
\end{itemize}

\begin{tcolorbox}[breakable,colback=teal!4,colframe=teal!35!black,title={Before you move on}]
What does \mr{हुकणे} mean? (\textquotedblleft{}To be missed; to miss a chance\textquotedblright{}.) Say it once more.
\end{tcolorbox}
